%% kapa: exam=1 round=29 year=2018 period=1 form=A subject=real_estate subject_ko=부동산학원론 seq=20 number=100
%% source: Q-Net 감정평가사 기출문제 (cst00309) artlSeq=5202808 file=extracted/1008/5202808/2018년 제29회 감정평가사 1차시험 기출문제/2018년 제29회 감정평가사 1차시험 기출문제/1교시 A형.pdf
%% conversion: kordoc 4.24.51 -> tools/gen1.py (auto); review=proofread-sonnet; flags=none
\begin{question}{100}
시장실패 또는 정부의 시장개입에 관한 설명으로 옳지 않은 것은?
\begin{choices}
  \item 외부효과는 시장실패의 원인이 된다.
  \item 소비의 비경합성과 비배제성을 수반하는 공공재는 시장실패의 원인이 된다.
  \item 정보의 비대칭성은 시장실패의 원인이 아니다.
  \item 시장가격에 임의로 영향을 미칠 수 있는 독과점 공급자의 존재는 시장실패의 원인이 된다.
  \item 시장실패의 문제를 해결하기 위하여 정부는 시장에 개입할 수 있다.
\end{choices}
\end{question}
\answer{3}
